% ch-convergence.tex -- numerical method: the correction floor, the convergence tests, termination
\chapter{Convergence and the correction floor}
\label{ch:convergence}

This is the most practically important chapter in the manual. Almost every
report of ``Yuzen is unbearably slow'' or ``the run stopped moving'' traces back to
one of two things: a convergence test that cannot be satisfied, which collapses
the time step, or a steady-state test that is satisfied far too early, which
stops the run while it is still going somewhere. Both come from the same
misunderstanding --- judging an iteration or a run by an absolute change ---
and both are avoidable once you know why the Newton correction has a floor it
cannot fall below.

Four keys decide when the Newton iteration of a step stops, all of them in the
block \key{Numerics.HeightEquation} (\cref{sec:ref-numerics:height}): they are
\key{NewtonEpsilon}, \key{NewtonEpsilonRel}, \key{NewtonMaxIterations} and \key{NewtonMinIterations}. Two
more decide when the run as a whole stops, in the block \key{Numerics}
(\cref{sec:ref-numerics:numerics}): \key{SteadyStateCriterion} and
\key{SteadyStateSteps}.

\section{The floor of the correction}
\label{sec:convergence:floor}

Newton's method converges quadratically, and for the first few iterations the
correction does exactly what that promises. Then it stops falling, and from
there on it merely scatters about one value. On the spreading film of
\cref{tab:floor} at \key{Nx = 2000}:

\begin{plain}
iteration 0 : 6.231e-02      <- quadratic convergence, falling cleanly
iteration 1 : 1.899e-03
iteration 2 : 1.494e-05
iteration 3 : 1.037e-09
iteration 4 : 1.931e-16      <- this is the "floor"
iteration 5 : 2.220e-16
iteration 6 : 2.189e-16      <- from here it wanders between 1.5e-16 and 2.2e-16
\end{plain}

The floor is \textbf{cancellation in the assembly of the right-hand side}, and
it has nothing to do with the accuracy of the linear solver. The right-hand
side is the residual of \cref{eq:nk-update},

\begin{equation}\label{eq:residual}
  \mathrm{Rhs} = (\varphi_{\mathrm{old}} - \varphi_{\mathrm{guess}})
    - \Delta t\left[\theta\,A\cdot\varphi_{\mathrm{guess}}
      + (1-\theta)\,A\cdot\varphi_{\mathrm{old}} + B\right],
\end{equation}

and $A$ is a fourth-order spatial difference, so its coefficients are of order
$1/\Delta x^{4}$: at \key{Nx = 2000} --- a non-dimensional
$\Delta x = 1.72\times10^{-4}$ --- the row sums reach $\sim 6\times10^{15}$.
Formed as the product of those entries with $\varphi$, $A\cdot\varphi$ is an
answer of about $10^{5}$ obtained by canceling terms of that magnitude
against one another, with an absolute error of
$\epsmach \times \lVert A\rVert \times \lvert\varphi\rvert \sim 1$ that
\cref{eq:residual} multiplies by $\Delta t$. Since the inverse of
$I + \theta\,\Delta t\,J$ is a contraction, the floor of such a product is
roughly

\begin{equation}\label{eq:floor}
  \epsmach \times \frac{\Delta t}{\dtexp},
  \qquad
  \dtexp \sim \frac{3\mu\,\Delta x^{4}}{H^{3}},
\end{equation}

proportional to how far beyond the explicit stability limit the step goes.
The well-mixed thickness equation does not form its stiff terms that way. The
capillary pressure and the vertical gravity give their $A\cdot\varphi$ face by
face (\cref{sec:time:apply}): the flux of a face is its mobility times a
difference of the film, and the differences are taken first, between
neighboring values. What cancels is then the variation of the film over the
stencil rather than the film, and on the films of \cref{tab:floor} the floor
is that of $\varphi$ itself, a few $\epsmach$, whatever the mesh and the step.
\Cref{eq:floor} remains the bound for what is formed as a product --- the
in-plane body forces, the Marangoni stress and the disjoining pressure, whose
entries are of order $1/\Delta x$ or $1/\Delta x^{2}$ --- and for a film whose
variation over a stencil is comparable to the film itself. Set
\env{YUZEN\_DEBUG\_FLOOR} to \val{1} and the floor is printed.

\begin{yzntable}{Measured correction floor, and the bound $\epsmach\,\Delta t/\dtexp$ of \cref{eq:floor} for $A\cdot\varphi$ formed as a product\label{tab:floor}}{@{}L{0.24\textwidth} d d d d d@{}}
\textsf{\bfseries Case} & \textsf{\bfseries $\Delta x$ [\si{\micro\metre}]} &
\textsf{\bfseries $\Delta t$} & \textsf{\bfseries $\Delta t/\dtexp$} &
\textsf{\bfseries bound} & \textsf{\bfseries measured} \\ \midrule
spreading, Nx = 500  & 1.2  & 3.21e-9  & 4.8e3 & 1.1e-12 & 1.1e-16 \\
spreading, Nx = 2000 & 0.3  & 3.21e-9  & 1.2e6 & 2.7e-10 & 1.8e-16 \\
spreading, Nx = 2000 & 0.3  & 3.21e-11 & 1.2e4 & 2.7e-12 & 1.1e-16 \\
flat film, \file{leveling\_1d\_reg} & 3.33 & 3.21e-8 & 8.1e2 & 1.8e-13 & 5.5e-17 \\
\end{yzntable}

The $\Delta t$ column is non-dimensional. The spreading film is
\file{stress\_wetting} at the \key{Nx} given, run with the adaptive step for
\SI{1}{\milli\second} (\SI{0.1}{\milli\second} at \key{Nx = 2000}) and then
stepped at the $\Delta t$ of the table, \key{NewtonEpsilon} and \key{NewtonEpsilonRel}
out of the way; the floor is the median of the corrections from the fifth
iteration on. Two conclusions follow.

\begin{description}
  \item[No particular value is universal.] On these films the floor does not
    grow as $\Delta x^{-4}\Delta t$, but a term formed as a product, a film
    steep over a few cells or another case can put it higher, as high as
    \cref{eq:floor}. Measure it rather than assume it.
  \item[Machine precision is the limit.] The correction does not fall below
    $\epsmach\,\lvert\varphi\rvert$: the updated film cannot be represented
    more finely. The linear solver is not what sets it --- the measured
    \code{bcg\_resid} is $\sim 10^{-25}$ --- and tightening its tolerance does
    not move the floor at all.
\end{description}

\section{What an absolute-only test does}
\label{sec:convergence:absolute}

\begin{warn}[Do not use an absolute-only convergence test]
  Where the floor is proportional to $\Delta t$ (\cref{eq:floor}), an
  absolute threshold below it creates the relation ``make $\Delta t$ smaller
  and the floor falls until the test is met''. That relation closes into a
  feedback loop with the time-step control, and the run locks onto the step
  at which the two just meet. A threshold below $\epsmach\,\lvert\varphi\rvert$
  is never met at all.
\end{warn}

The loop runs as follows.

\begin{enumerate}
  \item A \key{NewtonEpsilon} below the floor can never be satisfied.
  \item The iteration has effectively converged after about five iterations,
    but it spins on to \key{NewtonMaxIterations}.
  \item Having hit the iteration cap, it reports a convergence failure.
  \item The caller, \code{TimeAdvance}, sees the failure and \textbf{halves the
    time step}.
  \item A smaller step lowers the floor, so the run settles onto the step size
    that just reaches \key{NewtonEpsilon} and oscillates there --- \code{dt-} and
    \code{dt+} alternating in the log.
\end{enumerate}

The finer the mesh, the higher a floor of the kind of \cref{eq:floor} and the
smaller the step the run locks onto: on a mesh that resolves a contact line
the loop can cost orders of magnitude in the physical time reached.

\section{The two tests that avoid the floor}
\label{sec:convergence:tests}

\subsection{The relative test, \key{NewtonEpsilonRel}}

The iteration is taken to have converged once the correction has fallen to
\key{NewtonEpsilonRel} times \textbf{the first correction of that time step};
the default is \val{1e-6}. The statement it makes is ``the non-linearity of this
time step has been solved to six digits'', which is a statement about the step
and nothing else --- it is independent of the scale of the solution, of the
size of the time step and of the mesh. The floor and $\Delta t$ appear in the
numerator and the denominator alike and cancel, which is exactly what breaks
the feedback loop of \cref{sec:convergence:absolute}. The iteration starts
from the film extrapolated from the last two steps (\cref{sec:time:newton}),
so its first correction is what the extrapolation missed, and the test is
relative to that: taken relative to how far the extrapolation moved the film
instead, the iteration stopped a pass earlier, and the volume of a nozzle film
thickened a hundredfold drifted to $|\mathrm{VolErr}| = 5\times10^{-11}$
(\file{tools/\bk analytic.py} \code{nozzle\_balance}); relative to the first
correction it stays at $7\times10^{-14}$.

On the history at the start of this chapter the threshold is
$6.231\times10^{-2} \times 10^{-6} = 6.23\times10^{-8}$, so the iteration exits
at iteration 3 --- before it has reached the floor at all. Setting
\key{NewtonEpsilonRel} to \val{0} or below switches the relative test off and
leaves the absolute test alone.

\subsection{The stagnation test}
\label{sec:convergence:stagnation}

The second test looks at \textbf{how many iterations in a row fail to improve
the best value so far}. An iteration that does not get below $0.9$ times the
smallest correction seen so far counts as one failure, and more than five
failures stops the loop.

The detail that makes it work is the use of the best value rather than the
change from one iteration to the next. Measuring the change ratio
$\lvert\mathrm{resid} - \mathrm{resid_{old}}\rvert/\mathrm{resid_{old}}$
instead does not fire at all in the situation it is meant to catch: on the
floor the correction scatters by tens of percent, and by several times where
the floor is of the kind of \cref{eq:floor}, from iteration to iteration,
which any change-ratio test reads as healthy progress.

When the test fires, the outcome depends on where the correction has got to.
The iteration has converged as far as it can, and the step is accepted
silently, if any of these holds:

\begin{itemize}
  \item the absolute test is met;
  \item the correction has fallen four or more decades below the first
    correction of the step: the non-linearity of the step is resolved to
    four digits;
  \item the correction is within $1000\,\epsmach\,\lvert\varphi\rvert$, with
    $\lvert\varphi\rvert$ the largest film value: the film cannot be updated
    more finely than $\epsmach\,\lvert\varphi\rvert$, and the floors of
    \cref{tab:floor} are about that. This catches a step whose first
    correction was already close to the floor, near a steady state, without
    having to lower \key{NewtonEpsilon}.
\end{itemize}

Otherwise the iteration has stopped \textbf{far from convergence}, and the
step has failed, like one whose linear solve does not converge or that runs
out of \key{NewtonMaxIterations}. With \key{Numerics.TimeStepControl.Enable =
true} it is cut by \key{ChangeFactor} and solved again, the first such step is
reported in \file{Warnings.txt} and all of them are counted in
\file{JobInfo.txt} (\code{Steps failed: the Newton correction stagnated},
\cref{sec:output:jobinfo}); a cut below \key{MinDt} aborts the run, as any
other (\cref{sec:time:dt}). At a fixed step there is no shorter step to retry
it at, so the run stops and names the stagnation
(\cref{sec:ref-numerics:tsc}).

At the default tolerances every stagnation in the shipped cases and the
checks that ship with them is of the second kind, and far into it: on the
contact-line stress case (\file{stress\_wetting}, the first step) the
correction stopped at 0.068 to 25 times the first correction, on the thick nozzle film
with liquid supplied at its own concentration (\code{nozzle\_supplied},
\cref{sec:verification:analytic}) at $8\times10^{-4}$ to 1.9 times; no
stagnation fell between the relative test ($10^{-6}$) and $8\times10^{-4}$.
Kept, such a step is not a solution of the step: the nozzle film kept its
residual as a volume error of $2.2\times10^{-8}$, and the contact-line film
run at a fixed step went negative and ended as diverged. Cut, the nozzle film
is at rounding ($2\times10^{-13}$).

The four decades are deliberately not tied to \key{NewtonEpsilonRel}. With
\key{NewtonEpsilonRel} at $10^{-11}$ (and \key{NewtonEpsilon} at $10^{-16}$) the same nozzle film levels off at
$10^{-11}$ to $10^{-4}$ of the first correction in two steps out of three ---
the limited upwinding at Courant numbers up to 150 is not smooth --- and
those steps are accepted; failing them would chase the relative test the way
the loop of \cref{sec:convergence:absolute} chases an absolute one. The run
cuts 33 steps that stagnated above $10^{-4}$, takes 3358 steps and ends
with a \code{VolErr} of $1.5\times10^{-9}$, against $2\times10^{-13}$ at
the default: what the accepted steps left. A tighter \key{NewtonEpsilonRel}
buys nothing on a film whose iteration cannot reach it.

\section{What to look at when it does not converge}
\label{sec:convergence:diagnosis}

\begin{enumerate}
  \item Set \env{YUZEN\_DEBUG\_RESIDUAL} and find out whether it is the
    thickness or the concentration that is holding the iteration up
    (\cref{sec:time:coupling}).
  \item Set \env{YUZEN\_DEBUG\_FLOOR} and find out where the correction stops
    falling.
  \item If the floor sits above \key{NewtonEpsilon}, \textbf{rely on the relative
    test rather than lowering \key{NewtonEpsilon}}. Lowering the absolute threshold
    is the maneuver of \cref{sec:convergence:absolute} run in reverse: it
    makes the test harder, which shrinks the step, which lowers the floor, and
    the run ends up slower rather than more accurate.
  \item If the warning about the correction stagnating does appear --- look in
    \file{Warnings.txt} for \code{correction stagnated}
    (\cref{sec:output:warnings}), and in \file{JobInfo.txt} for how many
    steps were cut for it ---
    then the iteration genuinely is not converging and the floor is not the
    explanation. Look for a non-smooth ingredient --- upwind switching, or the
    clamping applied to evaporation and to the concentration --- which can
    prevent the quadratic behavior the test assumes, or a step far too long
    for the state, as at the start of \file{stress\_wetting}.
    \env{YUZEN\_DEBUG\_RESIDUAL} prints a line for every stagnation, with
    the floor it was judged against.
\end{enumerate}

\section{Termination}
\label{sec:convergence:termination}

At the end of every step \code{CheckTermination()} applies three tests in
order.

\begin{enumerate}
  \item \textbf{The specified time, or the step count, has been reached}: the
    run has passed \key{Numerics.EndTime}, or it has taken
    \key{Numerics.MaxSteps} steps. The two are reported separately, as
    \code{Time reached to the EndTime.} and
    \code{Step count reached Numerics.MaxSteps.}, so a run that stops early
    does not have to be diagnosed by counting dots. \key{MaxSteps} defaults to
    \val{0}, no limit, so unless it is set the end time is what stops the
    run; see \cref{sec:ref-numerics:numerics}.
  \item \textbf{The film has dried out}: the concentration selected by
    \key{Physics.Evaporation.EndConcType} --- \val{"mean"}, \val{"max"} or
    \val{"min"} --- has exceeded \key{EndConcValue}. The run
    prints \code{Dry out.}
  \item \textbf{The solution has stopped changing}, which is the subject of the
    rest of this section.
\end{enumerate}

\begin{warn}[Do not judge steadiness on the change per step]
  $\max\lvert h_{\mathrm{new}} - h_{\mathrm{old}}\rvert /
   \max\lvert h_{\mathrm{new}}\rvert$ becomes small simply because the time
  step is small. A calculation that resolves the contact line at
  $\Delta t \sim 10^{-11}$ would be declared ``steady'' within a dozen steps of
  the start and stop without complaint. The more finely the case is resolved
  the sooner this happens, which is what makes it hard to notice.
\end{warn}

What the test actually uses is the relative change \textbf{per unit time},

\begin{equation}\label{eq:changerate}
  \mathrm{ChangeRate} =
    \frac{\max\lvert h_{\mathrm{new}} - h_{\mathrm{old}}\rvert}
         {\max\lvert h_{\mathrm{new}}\rvert \times \Delta t}
  \quad [\si{\per\second}],
\end{equation}

and the criterion applied to \cref{eq:changerate} is

\begin{equation}\label{eq:steady}
  \mathrm{ChangeRate} \times (\mathrm{EndTime} - \mathrm{StartTime})
    < \mathrm{SteadyStateCriterion}.
\end{equation}

The left-hand side of \cref{eq:steady} reads ``how much the shape would
change if it went on changing at the present rate for the whole span of the
run, from its start ($0$, or \key{Restart.RestartTime}) to \key{EndTime}''. It is
dimensionless and depends neither on the time step nor on the time
scale $T_0$. The default \key{SteadyStateCriterion} of \val{1e-3} therefore
means ``at the present rate, the whole span of the run would move it by only
\SI{0.1}{\percent}''.

Once \cref{eq:steady} has held for more than \key{SteadyStateSteps} consecutive
steps the run prints \code{Steady state.} and stops. The first five steps are
not judged at all --- \code{ChangeRate} is reported as 1 for them --- and
setting \key{SteadyStateCriterion} to \val{0} or below disables the test
entirely.

\code{ChangeRate} is written out as a column of \file{CalcInfo.dat}.
\textbf{Look at that column before you choose a threshold}, rather than
guessing one: it tells you directly what range the quantity takes in your case,
and whether the default is loose or tight for it.
